% sections/method.tex
% Note: this paper uses \sfrac{a}{b} from xfrac, NOT \nicefrac{a}{b} (nicefrac.sty unavailable).
\section{Method}
\label{sec:method}

\subsection{Problem setup}
\label{sec:setup}

We consider a continuous-control Markov decision process with state $s \in
\mathcal{S} \subseteq \R^{n_s}$, action $a \in \mathcal{A} \subseteq \R^{n_a}$,
and deterministic transition $s_{t+1} = f(s_t, a_t)$ implemented by a
black-box physics engine $\mathcal{E}$ (no gradient through $\mathcal{E}$).
The agent is a stochastic policy $\pi_\theta(a \mid s)$ with parameters
$\theta \in \R^p$ (a multilayer perceptron).

The viability of a trajectory $\tau = (s_0, a_0, \dots, s_H)$ is measured by
$k$ scalar margins $m_1(s_t), \dots, m_k(s_t)$, each positive when the
constraint is satisfied (height above floor, body upright, no foot
penetration, etc.). For Walker2d-v4 ($n_s = 17$, $n_a = 6$) we use $k = 3$:
torso height, pitch angle, and minimum foot-clearance
(Section~\ref{sec:exp}). The training objective is the soft-minimum
(log-sum-exp) over the trajectory:
\begin{equation}
  \label{eq:soft-margin}
  J_\beta(\theta) \;=\;
  \mathbb{E}_{\tau \sim \pi_\theta}\!\left[
    -\frac{1}{\beta}\,\log\sum_{t,j} e^{-\beta\, m_j(s_t)}
  \right].
\end{equation}
This is the only training signal: no auxiliary task return, no shaping, no
exploration bonus.

\subsection{Truncated Student-$t$ mollifier policy}
\label{sec:mollifier}

The policy is a per-coordinate, truncated Student-$t$:
\begin{equation}
  \label{eq:mollifier}
  \pi_\theta(a \mid s) \;=\;
  \prod_{i=1}^{n_a}
  \frac{t_{\nu(s)}\!\bigl(a_i;\, \mu_i(s),\, \sigma_i(s)\bigr)}
       {Z_i\bigl(\mu_i(s), \sigma_i(s), \nu(s),\, \bar a_i\bigr)},
\end{equation}
where $(\mu(s), \log\sigma(s), \log\nu(s))$ are the three output heads of
the network, $|a_i| \le \bar a_i$ is the per-actuator torque bound, and
$Z_i$ is the truncation constant. The Student-$t$ density $t_\nu$ is
$C^\infty$ in $a$ and has full support on $[-\bar a_i, \bar a_i]$. The
degrees of freedom $\nu(s)$ are state-dependent: at $\nu \to \infty$ the
distribution recovers the Gaussian (peaked, exploitative); at small
$\nu$ the heavy tails enable exploratory escape from local plateaus
arising at contact-mode boundaries.

This mollifying property is the central technical reason a value function
is unnecessary; we prove smoothness of $J_\beta$ in $\theta$ in
Section~\ref{sec:theory}.

\subsection{Capacity-clamped network}
\label{sec:capacity}

Let $\pi_\theta$ be an $L$-layer ReLU MLP with weights
$\{W_\ell, b_\ell\}_{\ell=1}^L$. After every gradient step we project
each weight matrix back to a Frobenius ball:
\begin{equation}
  \label{eq:clamp}
  W_\ell \;\leftarrow\;
  \frac{W_\ell}{\max(1,\, \|W_\ell\|_F / c_\ell)},
\end{equation}
where the per-layer cap $c_\ell$ is determined from the
control-authority bounds of the task (Section~\ref{sec:hyper}). The
clamp plays the role of $L_2$ weight decay but is sharper: it enforces
a hard Lipschitz bound on each layer rather than a soft penalty in the
loss.

\subsection{Algorithm}
\label{sec:algorithm}

\begin{algorithm}[H]
\caption{Minimalist policy gradient (\textsc{MinPG}).}
\label{alg:minpg}
\begin{algorithmic}[1]
\Require Engine $\mathcal{E}$, env model (XML); margin functions
$\{m_j\}_{j=1}^k$; horizon $H$; eval count $N$;
network depth $L$, width $n$;
learning rate $\eta$; soft-min sharpness $\beta$;
capacity caps $\{c_\ell\}$.
\State Initialise $\theta = \{W_\ell, b_\ell\}_{\ell=1}^L$ randomly,
then apply clamp~\eqref{eq:clamp} per layer.
\For{epoch $= 1, 2, \dots$}
  \State Sample $N$ rollouts: for each, draw
  $a_t \sim \pi_\theta(\cdot \mid s_t)$, step $\mathcal{E}$, record
  $(s_t, a_t)_{t=0}^H$.
  \State For each rollout compute $J_\beta$ via~\eqref{eq:soft-margin}.
  \State Compute score-function gradient
  $\hat g \;=\; \frac{1}{N}\sum_n J_\beta^{(n)} \cdot
  \sum_{t=0}^H \nabla_\theta \log \pi_\theta(a_t^{(n)} \mid s_t^{(n)})$.
  \State Update $\theta \leftarrow \theta + \eta\, \hat g$.
  \State Re-apply clamp~\eqref{eq:clamp} to every $W_\ell$.
\EndFor
\end{algorithmic}
\end{algorithm}

\paragraph{What is absent.} \textsc{MinPG} has no value baseline, no
GAE, no replay buffer, no clipped surrogate, no entropy bonus, no
target network, no DAgger correction, no reward normalisation, and no
gradient through $\mathcal{E}$. The score function is computed through
$\pi_\theta$ only.

\subsection{Hyperparameter atlas}
\label{sec:hyper}

The seven scalars that govern \textsc{MinPG} split sharply into two
tiers.

\paragraph{Tier I (physics-determined; not free choices).}
$\eta$ -- the learning rate, bounded above by the contractivity
condition $\eta \ll 1/\gamma$ where $\gamma$ is the simulator's
intrinsic damping (read from the engine XML).
$H$ -- the horizon, set by the task: typically $500$--$2000$ steps at
the engine timestep $\Delta t$.
$\{c_\ell\}$ -- per-layer capacity caps, derived from the
torque/velocity envelope of the actuators.

\paragraph{Tier II (user choices; tune on validation rollouts).}
$L \in \{3, 4, 5\}$ -- depth.
$n \in \{128, 256, 512\}$ -- width.
$N \in \{64, \dots, 4096\}$ -- rollouts per gradient step.
$\beta \in \{1, \dots, 100\}$ -- soft-min sharpness; see
Remark~\ref{rem:annealing}.

\begin{remark}[Sharpness annealing]
\label{rem:annealing}
Small $\beta$ averages the gradient across all timesteps; large $\beta$
concentrates on the worst-margin step. We anneal $\beta$ upward with
training progress: explore broadly first, then sharpen the focus.
\end{remark}
